\section{A real tensor with uniform sign separation}
\label{sec:gaussian}
The dimension bound is now proved. We turn to the independent
splitting theorem that will supply absorption from finite nuclear
dimension. Its first ingredient is purely real and finite-dimensional:
one tensor must take both signs on every sufficiently separated product
of two-point sets. The following lemma obtains this simultaneously over
all such configurations. We equip real tensor powers with the Euclidean
tensor product inner product.

\begin{lemma}\label{lem:gaussian}
Fix \(k\geq2\) and \(0<r<1\). There are \(n\geq1\) and \(G\in(\R^k)^{\otimes n}\) such that, for every list of pairs of unit vectors \((u_j^0,u_j^1)_{j=1}^n\) satisfying
\[
 |\langle u_j^0,u_j^1\rangle|\leq r\qquad(1\leq j\leq n),
\]
the numbers
\[
 \big\langle G,u_1^{\alpha_1}\otimes\cdots\otimes u_n^{\alpha_n}\big\rangle,
 \qquad \alpha\in\{0,1\}^n,
\]
include a value greater than \(1/2\) and a value less than \(-1/2\).
\end{lemma}
\begin{proof}
Choose
\[
 0<a<\min\left\{\log2-\log(4/3)-\tfrac14\log3,\ -\tfrac14\log r\right\}.
\]
The first expression is positive. A binary Hamming ball of radius \(\lfloor n/4\rfloor\) has cardinality at most
\[
 \sum_{b=0}^{\lfloor n/4\rfloor}\binom nb
 \leq 3^{n/4}\sum_{b=0}^n\binom nb3^{-b}
 =3^{n/4}(4/3)^n.
\]
Gilbert's greedy deletion argument \cite[Theorem~1]{Gilbert}
therefore produces, for all sufficiently large \(n\), a set \(S_n\subset\{0,1\}^n\) of cardinality \(N=\lfloor e^{an}\rfloor\) whose distinct words have Hamming distance at least \(n/4\).

Let the coordinates of \(G\) be independent standard real Gaussian variables. Fix an admissible list and set
\[
 X_\alpha=\langle G,u_1^{\alpha_1}\otimes\cdots\otimes u_n^{\alpha_n}\rangle
 \quad(\alpha\in S_n).
\]
These variables are jointly centered Gaussian. Their covariance matrix \(\Sigma\) has diagonal entries one and off-diagonal entries of absolute value at most \(r^{n/4}\). Consequently its eigenvalues lie in \([1-\lambda_n,1+\lambda_n]\), where
\[
 \lambda_n=N r^{n/4}\longrightarrow0.
\]
For large \(n\), the density of this Gaussian vector is bounded pointwise by
\[
 \left(\frac{1+\lambda_n}{1-\lambda_n}\right)^{N/2}
\]
times the density of \(N\) independent centered Gaussians of variance \(1+\lambda_n\). Indeed, \(\det\Sigma\geq(1-\lambda_n)^N\) and \(\Sigma^{-1}\geq(1+\lambda_n)^{-1}I\).

Writing \(\Phi\) for the standard normal distribution function gives
\begin{equation}\label{eq:prob}
 \Pr(X_\alpha\leq1\ \forall\alpha\in S_n)
 \leq\left[\sqrt{\frac{1+\lambda_n}{1-\lambda_n}}\,
 \Phi\big((1+\lambda_n)^{-1/2}\big)\right]^N
 \leq q^N
\end{equation}
for one \(q<1\) and every sufficiently large \(n\). The last inequality is uniform over admissible lists because the bracket tends to \(\Phi(1)<1\). By symmetry the same bound holds for the event \(X_\alpha\geq-1\) for all \(\alpha\).

Set \(\eta_n=(4n k^{n/2})^{-1}\). The compact space of admissible lists has an \(\eta_n\)-net, in the maximum Euclidean distance between corresponding vectors, of size at most
\[
 (C_k/\eta_n)^{2kn}.
\]
For example, mesh the coordinate cubes and select an admissible representative from each occupied box. The logarithm of this bound is \(O(n^2)\). A union bound using \eqref{eq:prob} shows that the probability that some net list lacks a value above \(1\) or below \(-1\) is at most
\[
 2(C_k/\eta_n)^{2kn}q^N\longrightarrow0.
\]
Also \(\mathbb E\|G\|^2=k^n\), so
\(\Pr(\|G\|\leq2k^{n/2})\geq3/4\).
For large \(n\), both the successful-net event and this norm bound therefore hold for some \(G\).

If an admissible list is approximated by a net list, telescoping the elementary tensors of unit vectors bounds the change of each evaluation by
\[
 \|G\|\,n\eta_n\leq\tfrac12.
\]
The two evaluations exceeding \(1\) and below \(-1\) on the net list retain values above \(1/2\) and below \(-1/2\), respectively. This proves the assertion for every admissible list.
\end{proof}

